\documentclass[11pt,a4paper]{article}
\usepackage[T1]{fontenc}
\usepackage{mathpazo}
\usepackage[margin=26mm,headheight=14pt]{geometry}
\usepackage{microtype,amsmath,amssymb,amsthm,mathtools}
\usepackage{booktabs,array,longtable,enumitem}
\usepackage{xcolor,xurl,hyperref,fancyhdr}
\definecolor{inkblue}{HTML}{193C60}
\hypersetup{colorlinks=true,linkcolor=inkblue,citecolor=inkblue,urlcolor=inkblue,
 pdftitle={Two reflections and exact ranks for level-four double polylogarithms},
 pdfauthor={OpenAI-assisted research contribution for ProveIt}}
\pagestyle{fancy}\fancyhf{}
\fancyhead[L]{\small Level-four double-shuffle ranks}
\fancyhead[R]{\small ProveIt research continuation}
\fancyfoot[C]{\thepage}
\setlength{\parskip}{3pt}
\setlength{\emergencystretch}{2em}
\numberwithin{equation}{section}
\newtheorem{theorem}{Theorem}[section]
\newtheorem{proposition}[theorem]{Proposition}
\newtheorem{corollary}[theorem]{Corollary}
\newtheorem{lemma}[theorem]{Lemma}
\theoremstyle{definition}\newtheorem{definition}[theorem]{Definition}
\theoremstyle{remark}\newtheorem{remark}[theorem]{Remark}
\newcommand{\Q}{\mathbb Q}
\newcommand{\Z}{\mathbb Z}
\newcommand{\ii}{\mathrm i}
\newcommand{\Li}{\operatorname{Li}}
\newcommand{\Imag}{\operatorname{Im}}
\newcommand{\rank}{\operatorname{rank}}
\newcommand{\Span}{\operatorname{span}}
\newcommand{\Row}{\operatorname{Row}}
\newcommand{\Hom}{\operatorname{Hom}}
\newcommand{\kerdim}{\operatorname{nullity}}
\newcommand{\rev}{\mathsf T}
\newcommand{\refl}{\mathsf R}
\newcommand{\Gset}{\mathcal G}
\newcommand{\Nset}{\mathcal N}
\newcommand{\V}{\mathcal V}
\newcommand{\R}{\mathcal R}
\newcommand{\Ltwo}{\log 2}
\newcommand{\pathtext}[1]{\nolinkurl{#1}}
\title{\vspace{-9mm}\textbf{Two reflections and exact ranks for\\level-four double polylogarithms}\\[5pt]
\large A proof of the manuscript's odd-weight rank conjecture,\\an explicit parity compiler, and uniform harmonic-sum obstructions}
\author{OpenAI-assisted research contribution prepared for\\Vladimir Reshetnikov's ProveIt project}
\date{9 October 2026, Pacific time}
\begin{document}
\maketitle
\begin{abstract}
The ProveIt manuscript records an all-odd-weight conjecture for a particular
rational matrix of level-four depth-two shuffle and stuffle product relations.
We prove both asserted ranks, not only their difference. If $A_w$ is the
matrix on its $6w-7$ convergent imaginary coordinates, then
$\rank A_w=5w-6$ for odd $w\ge3$, and $\rank A_w=6w-7$ for even $w\ge2$.
After deleting the Gaussian columns, the corresponding ranks are
$(9w-11)/2$ and $5w-6$. The proof converts the binomial rows into equations
for six homogeneous polynomials. Two independent polynomial reflections
then give an explicit integer basis of the entire odd-weight nullspace.
This yields a complete row-membership test and proves that the Gaussian-only
linear consequences are exactly the same-point shuffle consequences.

Keeping the single-value right sides gives a closed polynomial compiler for
all even-weight imaginary double values at fourth roots of unity. We extract
two compact mixed-color identities in every even weight. A second consequence
is a uniform obstruction: in every odd weight, the sum
$S_{2m}=\sum_{n\ge0}(-1)^nH_n/(2n+1)^{2m}$ introduces one additional formal
direction beyond the Gaussian doubles in this restricted linear system.
The proposed numerical evaluation of $S_4$ is neither proved nor refuted.
Exact certificates cover weights $2$--$41$, and 210 explicit even-weight
formulas through weight $12$ are checked against all 816 product rows.
All ranks are formal; no numerical period independence is asserted.
\end{abstract}
\vspace{3mm}
\noindent\textbf{Keywords.} Multiple polylogarithms; colored double zeta values;
double shuffle; Gaussian periods; exact linear algebra; parity.
\newpage
\tableofcontents
\newpage

\section{Research target, result, and scope}
\label{sec:scope}
The source for this continuation is the consolidated ProveIt manuscript at
commit \pathtext{9bc738d3be22b8586a24693f19fb2e2a50ecd1bf}. The target is the
conjecture labeled \pathtext{signed:conj:rank} in
\pathtext{chapters/05-signed-kernels.tex}; the underlying file has Git blob
\pathtext{e61ff32263722d7575ca0202fba4dd7dad534861}. That conjecture gives two
ranks for an explicitly specified family of rational matrices, with exact
experimental evidence through odd weight $31$ \cite{repo}.

This target is distinct from the manuscript's previously proved Gaussian
weight-six formulas, its earlier shuffle-rank theorem, and its distribution
rank results. The result below settles the remaining matrix conjecture and
explains its experimental pattern. The method also covers the even weights
and produces an explicit solution operator for the inhomogeneous equations.

\begin{theorem}[All-weight rank theorem]\label{thm:main}
Let $A_w$ be the level-four imaginary depth-two product matrix defined in
Section~\ref{sec:model}. Let $A_w^{\neg g}$ be obtained by deleting all columns
$\Imag\Li_{a,w-a}(\ii,1)$, $1\le a<w$. Then
\begin{align}
\rank A_w&=
\begin{cases}
5w-6,&w\ge3\text{ odd},\\
6w-7,&w\ge2\text{ even},
\end{cases}\label{eq:main-rank}\\
\rank A_w^{\neg g}&=
\begin{cases}
(9w-11)/2,&w\ge3\text{ odd},\\
5w-6,&w\ge2\text{ even}.
\end{cases}\label{eq:main-nong}
\end{align}
For odd $w$, an explicit integer basis of $\ker A_w$ has $w-1$ members,
split into two families of $(w-1)/2$ members. One family controls the Gaussian
coordinates; the other vanishes on every Gaussian coordinate.
\end{theorem}

The rank refers to coefficient rows after the known single-value right sides
have been suppressed. It is not a dimension of numerical periods. If the
evaluation map has further identities, those identities can reduce the
numerical span without changing $A_w$. Similarly, an obstruction in this
matrix system excludes only a rational linear combination of its stated
rows. It does not exclude argument transformations, elimination from higher
depth, or products of lower-weight relations followed by further reduction.

The even-weight reduction phenomenon belongs to the established parity theory
of multiple polylogarithms, notably Panzer's theorem \cite{panzer}. Our
contribution here is an explicit description of this particular linear
product module and its compiler; we do not claim a new general parity
principle. Standard cyclotomic relation systems are substantially broader
than this module \cite{zhao}. The literature search also located recent work
on a different, weight-two conjecture of Zhao \cite{bachmann}; it should not
be conflated with the manuscript conjecture proved here. No exhaustive claim
of worldwide priority is made for individual special-value identities.

\section{The exact product model}
\label{sec:model}
\subsection{Colors, indices, and convergent coordinates}
Our convention is
\begin{equation}
D_{a,b}(x,y)=\Li_{a,b}(x,y)
 =\sum_{n>m\ge1}\frac{x^ny^m}{n^a m^b},\qquad a,b\ge1.
\label{eq:series}
\end{equation}
Thus $a$ is the outer index. We use
\[
H_n^{(b)}=\sum_{j=1}^n j^{-b},\quad H_n=H_n^{(1)},\quad H_0^{(b)}=0,
\qquad \beta(k)=\sum_{j\ge0}\frac{(-1)^j}{(2j+1)^k}\quad(k\ge1),
\]
and $\zeta(k)=\sum_{n\ge1}n^{-k}$ for $k\ge2$.
The series for $\beta(1)$ is understood in its ordinary alternating sense.
Put
\[
x_{a,b}^{r,s}=\Imag D_{a,b}(\ii^r,\ii^s),\qquad r,s\in\Z/4\Z.
\]
Conjugation gives $x_{a,b}^{r,s}=-x_{a,b}^{-r,-s}$. When both colors are even,
the imaginary coordinate is zero. The six remaining color representatives are
\begin{equation}
(0,1),\ (1,0),\ (1,1),\ (1,2),\ (1,3),\ (2,1).
\label{eq:colors}
\end{equation}
For each representative, retain $a=1,\ldots,w-1$, except for
$x_{1,w-1}^{0,1}$. This last coordinate has a logarithmically divergent
imaginary part. There are consequently $6(w-1)-1=6w-7$ retained coordinates.
We order them lexicographically by $(a,b,r,s)$, with $b=w-a$.

For $a>1$, absolute convergence follows from
$H_{n-1}=O(\log n)$ and $H_{n-1}^{(b)}\le H_{n-1}$.
For $a=1$ and outer color different from $1$, summation by parts gives the
usual convergent boundary values. For example, when the inner color is $1$,
$H_{n-1}^{(b)}/n$ tends to zero and has summable eventual variation; when the
inner color is nontrivial, its partial sums converge to a finite single value
with an $O(n^{-b})$ remainder. The exceptional outer color $1$ and imaginary
inner color produces the removed logarithm. Purely real divergent expressions
are not assigned finite numerical values: their imaginary equations are
simply zero.

\subsection{The row vocabulary}
The stuffle product is
\begin{equation}
D_{p,q}(x,y)+D_{q,p}(y,x)
 =\Li_p(x)\Li_q(y)-\Li_w(xy),\qquad p+q=w.
\label{eq:stuffle}
\end{equation}
The indices \emph{and the colors} exchange together. The product shuffle is
\begin{align}
&\sum_{j=0}^{p-1}\binom{q+j-1}{j}
       D_{q+j,p-j}(y,x/y)\notag\\
&\qquad+\sum_{j=0}^{q-1}\binom{p+j-1}{j}
       D_{p+j,q-j}(x,y/x)
 =\Li_p(x)\Li_q(y).
\label{eq:shuffle}
\end{align}
The first identity follows by partitioning the product of two sums into
$n>m$, $m>n$, and $m=n$. The second is the iterated-integral shuffle, with the
binomial coefficients counting interleavings of the zero letters; its
coefficient form is also checked directly in Section~\ref{sec:polys}.
The boundary and regularization conventions agree with the source model
and standard regularized double-shuffle relations \cite{repo,zhao}.

For each $p+q=w$ and $r,s\in\Z/4\Z$, take imaginary parts with
$x=\ii^r$, $y=\ii^s$. If neither factor is $\Li_1(1)$, retain both coefficient
rows. If exactly one factor is $\Li_1(1)$, retain only the
\emph{stuffle-minus-shuffle} row. Its divergent double coordinate cancels and
its right side is $-\Imag\Li_w(xy)$. If both factors are divergent, the
imaginary equation is zero. Discard zero rows; retain duplicates if desired.
This is $A_w$. Repeated rows affect neither rank nor row membership.

For example, the implementation with repetitions has $92$ rows and $23$
columns at weight $5$, exactly matching the canonical manuscript model.
We do not add distribution relations, arbitrary changes of argument, or
higher-depth relations. Suppressing the single-value right sides gives the
homogeneous coefficient system used in the rank proof.

\subsection{A formal quotient and its evaluation meaning}
Let $\V_w=\Q^{6w-7}$, viewed as row vectors of formal double-coordinate
coefficients, and set
\[
\R_w=\Row(A_w),\qquad \mathcal Q_w=\V_w/\R_w.
\]
The quotient suppresses the single-value right sides of the specified rows.
It is deliberately not the algebra of all colored periods modulo all
products. A vector in $\R_w$ yields a valid reduction to a particular linear
combination of the displayed single-value right sides. The converse need not
hold for true numerical identities. This separation is essential when using
$\mathcal Q_w$ to reject an insufficient proof search.

\section{Completing the missing coordinate and encoding the rows}
\label{sec:polys}
Set $N=w-2$, and introduce six homogeneous polynomials of degree $N$:
\begin{equation}
\begin{array}{c|cccccc}
\text{color}&(0,1)&(1,0)&(1,1)&(1,2)&(1,3)&(2,1)\\ \hline
\text{polynomial}&A&B&C&D&E&F.
\end{array}
\label{eq:poly-names}
\end{equation}
For example,
$B(X,Y)=\sum_{a+b=w}v_{a,b}^{1,0}X^{a-1}Y^{b-1}$ for a candidate null vector
$v$. The coefficient absent from the matrix occurs in $A$.

\begin{lemma}[Unique homogeneous completion]\label{lem:completion}
Every vector in $\ker A_w$ extends uniquely to six complete polynomials
satisfying all homogeneous shuffle and stuffle equations. The missing
coefficient is assigned
\begin{equation}
v_{1,w-1}^{0,1}=-v_{w-1,1}^{1,0}.
\label{eq:completion}
\end{equation}
Restriction to the retained coordinates is an isomorphism back to $\ker A_w$.
\end{lemma}
\begin{proof}
Assigning \eqref{eq:completion} makes the missing stuffle equation zero.
Its shuffle equation is then zero because the retained difference equation
is zero. This handles the two orders of the divergent factor and their
conjugates. The absent coefficient occurs nowhere else: in a shuffle sum it
can have outer index $1$ only when $q=1,j=0$ or $p=1,j=0$, precisely the
divergent-factor cases. All other equations were already retained. Conversely,
a complete homogeneous solution satisfies every retained row. The missing
stuffle equation forces \eqref{eq:completion}, giving uniqueness. The argument
includes $w=2$; the equation with two divergent factors is purely real.
\end{proof}

For a homogeneous polynomial $P$, write $P^\rev(X,Y)=P(Y,X)$. The stuffle
equations reduce to
\begin{equation}
A=-B^\rev,\qquad C=-C^\rev,\qquad
D=-F^\rev,\qquad E=E^\rev.
\label{eq:stuff-polys}
\end{equation}
For instance, swapping the colors $(1,3)$ gives $(3,1)$, which is minus
$(1,3)$ under conjugation; this explains the positive symmetry of $E$.

\begin{lemma}[Polynomial form of the shuffle rows]\label{lem:shuffle-polys}
The remaining equations are exactly
\begin{align}
E(X+Y,X)+A(X+Y,Y)&=0,\label{eq:sh-A}\\
B(X+Y,X)+B(X+Y,Y)&=0,\label{eq:sh-B}\\
C(X+Y,Y)-F(X+Y,X)&=0,\label{eq:sh-C}\\
D(X+Y,Y)-D(X+Y,X)&=0.\label{eq:sh-D}
\end{align}
Together with \eqref{eq:stuff-polys}, these equations include all colors and
all weight splits, including the completed boundary equations.
\end{lemma}
\begin{proof}
For general color polynomials $P_{r,s}$, the polynomial whose coefficient of
$X^{p-1}Y^{q-1}$ is the shuffle row is
\[
P_{s,r-s}(X+Y,X)+P_{r,s-r}(X+Y,Y).
\]
Indeed, the term with indices $(a,b)=(q+j,p-j)$ in the first summand has the
coefficient $\binom{q+j-1}{q-1}=\binom{q+j-1}{j}$, exactly as in
\eqref{eq:shuffle}. The second summand gives the other binomial sum.
Use the color pairs $(0,1)$, $(1,1)$, $(1,2)$, and $(1,3)$ in turn. Their
polynomials are the four displayed left sides. The pairs $(1,0)$ and $(2,1)$
just exchange $X,Y$ in equations already obtained. Conjugate pairs multiply
an equation by $-1$; even-even pairs are zero. This exhausts the possibilities.
\end{proof}

The binomial matrices have now disappeared. What remains is a pair of
reflection problems for ordinary polynomials.

\section{The two reflection blocks}
\label{sec:blocks}
\subsection{The block carrying Gaussian coordinates}
From \eqref{eq:sh-A} and $A=-B^\rev$, the invertible change of variables
$(X,Y)\mapsto(X+Y,X)$ gives
\begin{equation}
E(X,Y)=B(X-Y,X).
\label{eq:E-from-B}
\end{equation}
Equation \eqref{eq:sh-B} becomes
\begin{equation}
B(X,Y)+B(X,X-Y)=0.
\label{eq:B-reflection}
\end{equation}
Let $\varepsilon=(-1)^N$. Homogeneity and \eqref{eq:B-reflection} imply
\begin{align*}
E(Y,X)&=B(Y-X,Y)
 =\varepsilon B(X-Y,-Y)\\
&=-\varepsilon B(X-Y,X)=-\varepsilon E(X,Y).
\end{align*}
Because $E=E^\rev$, it follows that $(1+\varepsilon)E=0$.
If $N$ is even, $E=0$, hence $B=0$ and $A=0$. If $N$ is odd, the symmetry
of $E$ imposes no additional condition.

When $N=2m-1$ is odd, the change of variables $U=X$, $V=2Y-X$ turns
\eqref{eq:B-reflection} into oddness in $V$. Consequently the complete
Gaussian block has the basis
\begin{equation}
B_j(X,Y)=X^{N-2j-1}(2Y-X)^{2j+1},\qquad 0\le j<m,
\label{eq:B-basis}
\end{equation}
with $A_j=-B_j^\rev$ and
\begin{equation}
E_j(X,Y)=(X-Y)^{N-2j-1}(X+Y)^{2j+1}.
\label{eq:E-basis}
\end{equation}
Its dimension is exactly $m$.

\subsection{The complementary block}
Equations \eqref{eq:sh-C} and \eqref{eq:stuff-polys} give
\begin{equation}
F(X,Y)=C(X,X-Y),\qquad D(X,Y)=-C(Y,Y-X),
\label{eq:DF-from-C}
\end{equation}
where $C=-C^\rev$. The last shuffle equation asks for
$D(X,Y)=D(X,X-Y)$. But
\begin{align*}
D(X,X-Y)&=-C(X-Y,-Y)\\
&=-\varepsilon C(Y-X,Y)
 =\varepsilon C(Y,Y-X)=-\varepsilon D(X,Y).
\end{align*}
Thus $(1+\varepsilon)D=0$. If $N$ is even, $D=0$ forces $C=F=0$. If $N$ is
odd, the only condition is $C=-C^\rev$. A basis is
\begin{equation}
C_j(X,Y)=X^jY^{N-j}-X^{N-j}Y^j,\qquad 0\le j<m.
\label{eq:C-basis}
\end{equation}
This block also has dimension $m$ and vanishes identically on the Gaussian
polynomial $B$.

\begin{theorem}[Complete nullspace normal form]\label{thm:nullspace}
In every even weight $w\ge2$, $\ker A_w=0$. In odd weight $w=2m+1$, set
$N=2m-1$. Every vector in $\ker A_w$ is obtained uniquely by choosing
\begin{align*}
B&\in\Span_\Q\{B_0,\ldots,B_{m-1}\},\\
C&\in\Span_\Q\{C_0,\ldots,C_{m-1}\},
\end{align*}
then setting
\begin{equation}
(A,B,C,D,E,F)
 =\bigl(-B^\rev,\ B,\ C,\ -C(Y,Y-X),\ B(X-Y,X),\ C(X,X-Y)\bigr),
\label{eq:normal-form}
\end{equation}
and dropping the missing coefficient of $A$.
All resulting basis vectors have integer coordinates.
\end{theorem}
\begin{proof}
The preceding two block computations prove necessity and show that the two
blocks have disjoint determining polynomials. Substitution into
\eqref{eq:stuff-polys}--\eqref{eq:sh-D} proves sufficiency. The bases
\eqref{eq:B-basis} and \eqref{eq:C-basis} are independent respectively by the
change to $(U,V)$ and by their disjoint antisymmetric monomial pairs.
Deleting the missing coefficient cannot destroy independence, because the
$B$ and $C$ polynomials are retained in full. Lemma~\ref{lem:completion}
finishes the argument.
\end{proof}

\subsection{Proof of both rank formulas}
For even $w$, the column count and zero nullity give
$\rank A_w=6w-7$. Every column subfamily is independent, so deleting the
$w-1$ Gaussian columns gives rank $5w-6$.

For odd $w=2m+1$, Theorem~\ref{thm:nullspace} gives
$\kerdim A_w=2m=w-1$. Thus $\rank A_w=(6w-7)-(w-1)=5w-6$.
A null vector of $A_w^{\neg g}$ extends by zero on the Gaussian columns to a
null vector of $A_w$ with $B=0$. Exactly the $C$ block remains, of dimension
$m$. Since $A_w^{\neg g}$ has $5w-6$ columns,
\[
\rank A_w^{\neg g}=5w-6-m=\frac{9w-11}{2}.
\]
This proves Theorem~\ref{thm:main}, including the second rank separately.

\section{Gaussian saturation and an exact membership algorithm}
\label{sec:membership}
Let $\Gset_w\subset\V_w$ be the subspace supported on Gaussian columns.
Projection of the row space onto the other columns gives
\begin{equation}
\dim(\R_w\cap\Gset_w)=\rank A_w-\rank A_w^{\neg g}.
\label{eq:rank-difference}
\end{equation}

\begin{corollary}[Gaussian saturation in the stated module]\label{cor:saturation}
For odd $w\ge3$, the Gaussian-supported row space has dimension $(w-1)/2$
and is exactly the space spanned by the same-point product shuffles
$\Li_p(\ii)\Li_{w-p}(\ii)$, $1\le p\le(w-1)/2$.
For even $w$, the Gaussian-supported row space has dimension $w-1$.
\end{corollary}
\begin{proof}
The dimensions follow from \eqref{eq:rank-difference}. The indicated
same-point rows belong to the Gaussian-supported space: their shuffle
arguments are $(\ii,1)$. To check independence without importing any period
claim, write their coefficient on $g_{a,w-a}$ as
\[
\binom{a-1}{p-1}+\binom{a-1}{w-p-1}.
\]
For odd $w=2m+1$, restrict to $1\le a,p\le m$. The second binomial term
vanishes, and the remaining Pascal block is triangular with diagonal one.
There are $m$ independent rows in a space of dimension $m$.
\end{proof}

This is stronger than merely knowing that the small shuffle system has rank
$m$: it proves that eliminating every other color from the larger system
cannot add a new Gaussian-only coefficient relation at odd weight.

Let $K_w$ be the integer matrix whose columns are the basis vectors in
Theorem~\ref{thm:nullspace}. Finite-dimensional duality gives the following
complete decision procedure.
\begin{corollary}[Membership and separating certificates]\label{cor:membership}
For a rational row vector $t$,
\begin{equation}
t\in\R_w\quad\Longleftrightarrow\quad tK_w=0.
\label{eq:membership}
\end{equation}
When the product is nonzero, any column $v$ with $tv\ne0$ is an exact
certificate $A_wv=0$, $tv\ne0$ against membership. In even weight every
coefficient vector belongs to $\R_w$.
\end{corollary}
\begin{proof}
The row space is the annihilator of the nullspace over $\Q$. The displayed
columns are a basis of that nullspace, not just sampled null vectors.
\end{proof}

\subsection{Coefficient size and construction cost}
The basis can be constructed in $O(w^2)$ integer arithmetic operations.
Dehomogenize with $X=1$ and write $t=Y$. In the first block,
\[
B_j(1,t)=(2t-1)^{2j+1},\qquad
E_j(1,t)=(1-t)^{N-2j-1}(1+t)^{2j+1}.
\]
Successive $B_j$ arise by multiplication by $(2t-1)^2$. Successive $E_j$
arise by two exact divisions by $1-t$, followed by multiplication by
$(1+t)^2$. Each operation uses $O(w)$ coefficient updates. The divisions
have zero remainder by the displayed factorization and are implemented by
prefix sums with an explicit remainder check. The $C$ block has
\[
F_j(1,t)=(1-t)^{N-j}-(1-t)^j,
\]
and $D_j$ is minus its reversed coefficient vector. Binomial recurrences
therefore generate this block in the same bound.

All coefficients have $O(w)$ bits: the $B_j$ coefficients are bounded in
absolute value by $3^N$, the $E_j$ coefficients by $2^N$, and the remaining
binomial coefficients by $2^{N+1}$. A membership test then uses $O(w^2)$
integer arithmetic operations. For an integer target of coefficient height
at most $2^H$, a conservative schoolbook bit bound is
$O(w^2(H+w+\log w)^2)$, including basis construction. Rational inputs can
first be cleared of denominators; the bit size of that step must be charged
to the input representation. These are bounds for the membership task,
not for discovering every analytic relation among periods.

\section{A uniform obstruction for the even-index harmonic sums}
\label{sec:obstruction}
Define
\[
S_p=\sum_{n\ge0}\frac{(-1)^nH_n}{(2n+1)^p},\qquad
u_{p,1}=\Imag D_{p,1}(\ii,-1),\qquad
g_{p,1}=\Imag D_{p,1}(\ii,1).
\]
For $p\ge2$, absolute convergence permits separating the odd outer terms.
For an outer index $2n+1$,
\[
H_{2n}+\sum_{m=1}^{2n}\frac{(-1)^m}{m}=H_n.
\]
Consequently
\begin{equation}
S_p=g_{p,1}+u_{p,1}.
\label{eq:S-color}
\end{equation}

\begin{theorem}[All-odd-weight obstruction]\label{thm:S-obstruction}
For every odd $w=2m+1\ge3$, no rational linear combination of the rows of
$A_w$ can express $S_{2m}$ as a rational linear combination of Gaussian
doubles and the associated single-value products. More precisely, in the
formal quotient $\mathcal Q_w$, adjoining $S_{2m}$ to the Gaussian-coordinate
span increases its dimension from $m$ to $m+1$.
\end{theorem}
\begin{proof}
Choose the complementary null vector with
\[
B=0,\qquad C=C_0=Y^N-X^N,\qquad N=2m-1.
\]
It annihilates every Gaussian coordinate. Its $D$ polynomial is
$D(X,Y)=-C(Y,Y-X)$, so
\[
D(1,0)=-C(0,-1)=1.
\]
Therefore this null vector pairs to one with $u_{2m,1}$, hence also with
$S_{2m}$ in \eqref{eq:S-color}. It pairs to zero with every vector in the
row space and with every Gaussian-supported vector. This excludes a
reduction using the stated rows. The Gaussian span in $\mathcal Q_w$ has
dimension $m$ by Corollary~\ref{cor:saturation}; one new vector outside that
span increases the dimension by exactly one.
\end{proof}

\subsection{An integer certificate at weight five}
The source $S_4$ target has coefficient vector
\begin{equation}
T=3g_{4,1}+3g_{3,2}+9g_{2,3}+7u_{4,1}.
\label{eq:S4-target}
\end{equation}
Our $C_0$ witness has the following nonzero coordinates:
\begin{center}
\begin{tabular}{@{}crcr@{}}
\toprule
$(a,b,r,s)$&value&$(a,b,r,s)$&value\\\midrule
$(1,4,1,1)$&$1$&$(1,4,2,1)$&$-1$\\
$(2,3,1,2)$&$3$&$(2,3,2,1)$&$3$\\
$(3,2,1,2)$&$-3$&$(3,2,2,1)$&$-3$\\
$(4,1,1,1)$&$-1$&$(4,1,1,2)$&$1$\\\bottomrule
\end{tabular}
\end{center}
Thus $A_5v=0$ and $Tv=7$. The package records the full $92\times23$ matrix,
all four nullspace columns, the target, and this witness. In contrast to a
single isolated obstruction, the proof gives the same certificate pattern
in every odd weight. The previously supplied rational witness remains valid;
this is a simpler complementary-block alternative, not a correction to it.

The manuscript's proposed numerical $S_4$ evaluation remains open in this
article. The obstruction is fully compatible with its truth. A successful
proof must bring in a relation outside this fixed linear row vocabulary, or
use algebraic operations not already represented by these rows.

\section{The inhomogeneous compiler and exact odd-weight coordinates}
\label{sec:affine}
We now restore the right sides. This is also an independent route from the
rank theorem to explicit identities, without a search for integer relations.
Set
\[
\ell_k(r)=\Li_k(\ii^r),\qquad \ell_1(0)=0
\]
where the last assignment is a formal regularization used only in the
completed equations. It is not the value of a convergent series. Define
\begin{align}
P_{r,s}(X,Y)&=\sum_{p+q=w}\Imag(\ell_p(r)\ell_q(s))X^{p-1}Y^{q-1},
\label{eq:P-def}\\
H_N(X,Y)&=\sum_{j=0}^{N}X^jY^{N-j},\notag\\
S_{01}&=P_{0,1}-\beta(w)H_N,\qquad
S_{12}=P_{1,2}+\beta(w)H_N.\label{eq:S-def}
\end{align}
The notation $P$ here denotes a known product polynomial, not an unknown
coordinate polynomial. Define $A,B,C,D,E,F$ as in
\eqref{eq:poly-names}, now with the actual imaginary double values as
coefficients, and complete the missing coefficient by
\begin{equation}
[X^0Y^N]A=-g_{w-1,1}-\beta(w).
\label{eq:affine-completion}
\end{equation}
This is simply the missing stuffle equation with $\ell_1(0)=0$. The same
completion argument as Lemma~\ref{lem:completion} makes the system equivalent
to the original finite rows and their retained cancellation equations.

The inhomogeneous stuffle equations are
\begin{equation}
A+B^\rev=S_{01},\quad C+C^\rev=P_{1,1},\quad
D+F^\rev=S_{12},\quad E-E^\rev=P_{1,3}.
\label{eq:affine-st}
\end{equation}
The four shuffle right sides, in the order
\eqref{eq:sh-A}--\eqref{eq:sh-D}, are
$P_{0,1},P_{1,1},P_{1,2},P_{1,3}$ respectively.

Introduce the known polynomials
\begin{align}
e(X,Y)&=P_{0,1}(Y,X-Y)-S_{01}(X,X-Y),\label{eq:e}\\
d(X,Y)&=S_{12}(X,Y)+P_{1,2}(X,Y-X),\label{eq:d}\\
b(X,Y)&=P_{1,1}(Y,X-Y),\label{eq:b}\\
h(X,Y)&=P_{1,3}(Y,Y-X)-e(Y,Y-X)+e(Y-X,Y).
\label{eq:h}
\end{align}

\begin{theorem}[Closed even-weight compiler]\label{thm:compiler}
For even $w\ge2$, all the imaginary double coefficients are obtained from
\begin{align}
B(X,Y)&=\frac{b(X,Y)+h(X,Y)}2,\label{eq:B-aff}\\
C(X,Y)&=\frac12\bigl(P_{1,3}(X,-Y)+P_{1,1}(Y,X)\notag\\
&\hspace{26mm}-d(X-Y,-Y)+d(X-Y,X)\bigr),\label{eq:C-aff}\\
A(X,Y)&=S_{01}(X,Y)-B(Y,X),\label{eq:A-aff}\\
E(X,Y)&=B(X-Y,X)+e(X,Y),\label{eq:E-aff}\\
F(X,Y)&=C(X,X-Y)-P_{1,2}(Y,X-Y),\label{eq:F-aff}\\
D(X,Y)&=-C(Y,Y-X)+d(X,Y).\label{eq:D-aff}
\end{align}
The artificial coefficient in \eqref{eq:affine-completion} is discarded.
Every retained value is consequently an explicit rational combination of
single values and products of two single values.
\end{theorem}
\begin{proof}
Eliminating $A$ from the first affine shuffle gives
$E(X,Y)=B(X-Y,X)+e(X,Y)$. The second shuffle gives
\[
B(X,Y)+B(X,X-Y)=b(X,Y).
\]
Substitute the formula for $E$ in $E-E^\rev=P_{1,3}$ and use the even
homogeneity degree $N$; this gives
\[
B(X,Y)-B(X,X-Y)=h(X,Y).
\]
Adding proves \eqref{eq:B-aff}.

Similarly, eliminating $F$ gives \eqref{eq:D-aff}. The last shuffle states
$D(u,u-v)-D(u,v)=P_{1,3}(v,u-v)$. As $N$ is even,
\[
C(u-v,-v)=C(v-u,v)=P_{1,1}(v-u,v)-C(v,v-u).
\]
Substitution therefore gives
\begin{align*}
2C(v,v-u)= {}&P_{1,3}(v,u-v)+P_{1,1}(v-u,v)\\
&-d(u,u-v)+d(u,v).
\end{align*}
Put $(X,Y)=(v,v-u)$ to obtain \eqref{eq:C-aff}. The other formulas have
already been derived. Actual values satisfy the original equations, so
these necessary formulas evaluate them. Alternatively, the zero homogeneous
kernel proves uniqueness of the solution. No unproved independence of the
single values is needed.
\end{proof}

\begin{corollary}[Exact coordinates in the odd weights]\label{cor:odd-affine}
For odd $w=2m+1$, define
\[
B^\circ=B-\tfrac12 b,\qquad C^\circ=C-\tfrac12 P_{1,1}.
\]
Then $B^\circ(X,Y)=-B^\circ(X,X-Y)$ and
$C^\circ(X,Y)=-C^\circ(Y,X)$. They have unique expansions in the bases
\eqref{eq:B-basis} and \eqref{eq:C-basis}. Every other color polynomial is
recovered by \eqref{eq:A-aff} and \eqref{eq:E-aff}--\eqref{eq:D-aff}.
Thus precisely $m$ Gaussian and $m$ complementary coordinates remain in
this product system.
\end{corollary}
\begin{proof}
The polynomial $P_{1,1}$ is symmetric. Hence $b(X,Y)$ is invariant under
$Y\mapsto X-Y$. Subtract half of the two affine symmetrization equations
for $B$ and $C$. The homogeneous basis theorem then applies. The recovery
formulas did not require an assumption on the parity of $N$.
\end{proof}

For explicit scalar formulas, the single values used by the compiler are
\begin{align*}
\ell_1(1)&=-\tfrac12\Ltwo+\tfrac{\pi\ii}4,&\ell_1(2)&=-\Ltwo,\\
\ell_n(0)&=\zeta(n),&\ell_n(2)&=(2^{1-n}-1)\zeta(n),\quad n\ge2,\\
\ell_n(1)&=2^{-n}(2^{1-n}-1)\zeta(n)+\ii\beta(n),&
\ell_n(3)&=\overline{\ell_n(1)}.
\end{align*}
Even zeta values and odd beta values are then replaced by their standard
Bernoulli- and Euler-number expressions in powers of $\pi$. The implementation
keeps even beta values, odd zeta values, $\pi$, and $\log2$ as formal real
atoms. Testing equality of the resulting polynomial expressions does not
presume that their evaluated atoms are algebraically independent.

\section{Two all-weight mixed-color identities}
\label{sec:families}
The polynomial compiler has particularly simple endpoint coefficients.
Define
\[
c_p=\Imag D_{p,1}(\ii,\ii),\qquad
e_p=\Imag D_{p,1}(\ii,-\ii).
\]
These are different color patterns from $g_{p,1}=\Imag D_{p,1}(\ii,1)$.

\begin{theorem}[Mixed leading-index families]\label{thm:families}
For every integer $m\ge1$,
\begin{align}
c_{2m-1}={}&(m-1)\beta(2m)-\frac{\log2}{2}\beta(2m-1)\notag\\
&-\sum_{k=1}^{m-1}(1-2^{1-2k})\zeta(2k)\beta(2m-2k),
\label{eq:C-family}\\
e_{2m-1}={}&m\beta(2m)-\frac{\log2}{2}\beta(2m-1)\notag\\
&-\sum_{k=1}^{m-1}\zeta(2k)\beta(2m-2k).
\label{eq:E-family}
\end{align}
In particular,
\begin{equation}
e_{2m-1}-c_{2m-1}
=\beta(2m)-\sum_{k=1}^{m-1}2^{1-2k}\zeta(2k)\beta(2m-2k).
\label{eq:difference-family}
\end{equation}
\end{theorem}
\begin{proof}
Set $w=2m$, $N=w-2$. The desired coefficients are $C(1,0)$ and $E(1,0)$.
By \eqref{eq:C-aff},
\[
2C(1,0)=P_{1,3}(1,0)+P_{1,1}(0,1)-d(1,0)+d(1,1).
\]
The first two terms sum to $-(\log2)\beta(w-1)$. Definition \eqref{eq:d}
and $H_N(1,1)=N+1$ give
\[
d(1,1)-d(1,0)=P_{1,2}(1,1)-P_{1,2}(1,-1)+N\beta(w).
\]
In the difference of the two product polynomials, only even inner index
$q=2k$ remains. Its coefficient is
$2\beta(w-2k)\ell_{2k}(2)
=-2(1-2^{1-2k})\zeta(2k)\beta(w-2k)$.
Dividing by two proves \eqref{eq:C-family}.

For the second identity, \eqref{eq:B-aff} and \eqref{eq:E-aff} imply
\begin{align*}
2E(1,0)={}&P_{1,1}(1,0)+P_{1,3}(1,0)\\
&+e(1,0)+e(0,1).
\end{align*}
Again the first two terms give $-(\log2)\beta(w-1)$. Definition \eqref{eq:e}
yields
\[
e(1,0)+e(0,1)
=-P_{0,1}(1,1)+P_{0,1}(1,-1)+w\beta(w).
\]
Here $\ell_1(0)=0$ and $N$ is even. The product-polynomial difference is
$-2\sum_{k=1}^{m-1}\zeta(2k)\beta(w-2k)$. This proves
\eqref{eq:E-family}. Subtraction gives \eqref{eq:difference-family}.
\end{proof}

\subsection{Explicit examples}
With $G=\beta(2)$, the first weights are
\begin{align*}
c_1&=-\frac\pi8\log2,&e_1&=G-\frac\pi8\log2,\\
c_3&=\beta(4)-\frac{\pi^2G}{12}-\frac{\pi^3\log2}{64},&
e_3&=2\beta(4)-\frac{\pi^2G}{6}-\frac{\pi^3\log2}{64}.
\end{align*}
At weight six the formulas become
\begin{align}
\Imag\Li_{5,1}(\ii,\ii)
 &=2\beta(6)-\frac{\pi^2\beta(4)}{12}
   -\frac{7\pi^4G}{720}-\frac{5\pi^5\log2}{3072},\label{eq:c5}\\
\Imag\Li_{5,1}(\ii,-\ii)
 &=3\beta(6)-\frac{\pi^2\beta(4)}6
   -\frac{\pi^4G}{90}-\frac{5\pi^5\log2}{3072}.\label{eq:e5}
\end{align}
A non-endpoint example from the full compiler is
\begin{align}
\Imag\Li_{4,2}(\ii,-1)
={}&-\frac92\beta(6)-\frac{\pi^2\beta(4)}{48}
 +\frac{3\pi^3\zeta(3)}{64}+\frac{465\pi\zeta(5)}{512}.
\label{eq:d42}
\end{align}
The package contains all $210$ retained-coordinate identities in even weights
$2,4,6,8,10,12$. It separately checks both leading-index families through
weight $16$. These formulas are proved specializations, not PSLQ guesses.
Their identification as consequences of known product and parity structures
is preferable to an unsupported novelty claim for each individual number.

\section{Verification, provenance, and editorial corrections}
\label{sec:verification}
\subsection{Exact verification layers}
The core implementation constructs the coefficient matrix directly from
\eqref{eq:stuffle}--\eqref{eq:shuffle}, independently of the polynomial proof.
A second routine constructs the nullspace basis using coefficient recurrences,
not Gaussian elimination. For each tested weight, it verifies $A_wK_w=0$
as an exact integer matrix identity.

For rank lower bounds, a deterministic computation reduces the integer
matrices modulo the prime $65521$. A nonzero minor modulo that prime is a
nonzero rational minor. The explicit independent integer nullspace supplies
the matching rational upper bound. Thus the resulting rank certificates are
exact over $\Q$; they are not estimates of numerical rank and not a
probabilistic extrapolation from random primes. The same method handles the
non-Gaussian submatrix using the $C$ block. The central weight-five ranks are
also recomputed separately by rational elimination.

The default verification receipt records:
\begin{center}
\begin{tabular}{@{}p{.76\linewidth}r@{}}
\toprule
Verification scope&Count\\\midrule
All-weight matrix and deleted-column rank certificates, $2\le w\le41$&40\\
Odd-weight integer nullspace bases and uniform $S$ witnesses&20\\
Explicit even-weight coordinate formulas, $2\le w\le12$&210\\
Inhomogeneous product-row polynomial residuals&816\\
Odd-weight affine compatibility checks, $3\le w\le13$&6\\
Gaussian shuffle membership checks through weight $41$&420\\
Mixed leading-index coefficient checks through weight $16$&16\\
Independent floating-point direct-series diagnostics&54\\\bottomrule
\end{tabular}
\end{center}
The first seven rows are exact arithmetic checks. The last row is a separate
diagnostic and is not counted as a proof or an interval enclosure. The proof
of the infinite family is Sections~\ref{sec:polys}--\ref{sec:blocks}, not the
finite extent of the verification run.

\begin{samepage}
Selected exact rank results are
\begin{center}
\begin{tabular}{@{}rrrrr@{}}
\toprule
$w$&columns&$\rank A_w$&$\rank A_w^{\neg g}$&$\kerdim A_w$\\\midrule
2&5&5&4&0\\
3&11&9&8&2\\
4&17&17&14&0\\
5&23&19&17&4\\
6&29&29&24&0\\
7&35&29&26&6\\
15&83&69&62&14\\
31&179&149&134&30\\
40&233&233&194&0\\
41&239&199&179&40\\\bottomrule
\end{tabular}
\end{center}
\end{samepage}

\subsection{Independent direct-series diagnostics}
The numerical script does not use the shuffle rows to evaluate its left
sides. It sums \eqref{eq:series} directly with $4096$ outer terms at $70$
decimal working digits for $54$ cases with outer index at least $3$ in
weights $4,6,8$. The absolute mathematical truncation bound used for
comparison is
\[
\sum_{n>M}\frac{H_{n-1}^{(b)}}{n^a}\le
\begin{cases}
M^{1-a}\left(\dfrac{1+\log M}{a-1}+\dfrac1{(a-1)^2}\right),&b=1,\\[4pt]
\dfrac{b}{b-1}\dfrac{M^{1-a}}{a-1},&b\ge2.
\end{cases}
\]
It follows by the integral test and $H_{n-1}\le1+\log n$ or
$H_{n-1}^{(b)}\le b/(b-1)$. Every observed discrepancy is smaller than this
bound; the largest ratio is about $0.797$. Floating-point arithmetic in this
script is not outward rounded, so these comparisons are diagnostics rather
than fully certified intervals. Exact polynomial row residuals provide a
separate algebraic check on every delivered formula.

\subsection{A local correction to the source's logical wording}
The source states the two rank formulas and then calls the Gaussian-only
row-space dimension an equivalent statement. Rank-nullity gives
\eqref{eq:rank-difference}, so the pair of ranks implies the dimension of
the supported row space. The difference alone does not determine the two
ranks individually. The correct wording before an independent rank proof is
``In particular,'' not ``Equivalently.'' Theorem~\ref{thm:main} now proves
both ranks and removes the need for this overstatement.

The source's $92\times23$ obstruction matrix and its rational witness need
no correction. The previously repaired color-swapping rule also needs no
new correction: this article uses it essentially. This audit is focused on
the rank section and its algebraic dependencies; it is not a fresh,
exhaustive audit of all 228 manuscript pages.

\subsection{Integration and trust boundaries}
The supplied integration fragment replaces the last rank-conjecture
subsection in \pathtext{05-signed-kernels.tex}. It preserves the existing
labels \pathtext{signed:conj:rank} and
\pathtext{signed:eq:rank-conjecture} as compatibility aliases, but places
them on a theorem and a proved formula. The fragment includes a proof and
new labels under the prefix \pathtext{l4rank:}. Historical arrival sources
and receipts should remain unchanged as provenance.

The package includes the article source and compiled PDF, the matrix and
compiler routines, the exact verifier, a separate numerical script,
JSON certificates, an editorial audit, provenance metadata, and an optional
integration helper that refuses to operate on an unexpected Git blob.
The remote repository has not been modified. The ordinary mathematical
proofs have not been checked by Lean or another proof assistant, and the
article has not undergone independent peer review. These qualifications do
not turn the proved matrix theorem into a numerical conjecture; they specify
the verification boundary.

\section{Further research questions and next proof targets}
\label{sec:future}
\paragraph{The missing $S_4$ relation.}
Test each proposed additional law against the explicit complementary
$C_0$ witness. A law annihilating this witness cannot by itself remove the
obstruction. The source's suggested M\"obius transformation
$z\mapsto(1-z)/(1+z)$, or elimination from depth three and above, is a
concrete next source of candidate rows. The task is to derive an exact law
with a nonzero witness pairing and the correct single-value right side,
not merely a close numerical relation.

\paragraph{The complementary block in every odd weight.}
The polynomial $C$ now supplies explicit coordinates for the part missed by
Gaussian-only searches. Determine which additional relations act on these
coordinates and whether they have a uniform form. Since all mixed leading
harmonic sums already detect $C_0$, this is a systematic target rather than
a separate unstructured search at each weight.

\paragraph{Other cyclotomic levels.}
At level four the six conjugation classes split into two uncoupled reflection
blocks. At levels three, six, and eight, write the analogous color-polynomial
system and determine its orbit decomposition. The objective is a proven
nullspace parametrization, not a fit to a dimension sequence. Modular-form
and period-polynomial methods at general level offer related structures
\cite{hirose}, but their relation vocabulary must be matched carefully to the
one used here.

\paragraph{The real-part sector.}
The real sector retains even-even color coordinates and has different
conjugation signs. Its boundary completion also includes different divergent
components. A complete analogous analysis would show precisely where the
two-reflection simplification survives and where extra blocks arise.

\paragraph{Integral structure and Smith normal forms.}
The displayed basis is integral, but this article does not claim it is a
basis of the full integral kernel lattice. Determine the lattice index and
the Smith normal form of independent row submatrices. These invariants could
control unavoidable denominators in rational reduction certificates and
explain which small primes are exceptional.

\paragraph{A faster constructive positive certificate.}
The membership test immediately returns a separating vector when the answer
is negative. When the answer is positive, extend the reflection eliminations
to output an explicit sparse combination of the original matrix rows for the
target. The even-weight compiler is already constructive, but a uniformly
sparse, row-indexed certificate format would improve independent verification
and integration into a proof assistant.

\paragraph{Formal verification.}
The purely algebraic core is a feasible first formalization target: binomial
coefficient extraction, completion of one missing coordinate, homogeneous
polynomial substitutions, and dimensions of two eigenspaces. The analytic
shuffle and regularized boundary laws can initially be explicit hypotheses,
then discharged in a separate development. This separates the finite algebra
from the analytic theory rather than mixing their trusted assumptions.

\paragraph{Arithmetic identities beyond the formal quotient.}
The $w-1$ formal odd-weight coordinates need not be independent numerical
periods. Study their images under known motivic or analytic constructions,
stating every additional relation or period-map assumption explicitly.
The rank theorem gives a reliable coordinate system for this question, not
an arithmetic independence result.

\section{Conclusion}
The all-odd-weight matrix conjecture is settled by a complete, uniform
nullspace description. Its two ranks arise from two independent degree
$w-2$ polynomial eigenspaces, and the same calculation makes the even-weight
matrix injective. This resolves the source conjecture, establishes exact
Gaussian saturation for the specified product system, and provides a
quadratic-arithmetic membership test with explicit separating certificates.

Retaining the right sides gives a full even-weight compiler and the mixed
families \eqref{eq:C-family}--\eqref{eq:E-family}. The complementary block
explains the $S_4$ proof obstruction in every odd weight. The next genuinely
new step for $S_4$ is now sharply identified: find an analytically valid law
outside this module that acts nontrivially on its complementary coordinates.

\appendix
\section{Reproduction commands and data contracts}
\label{app:reproduction}
From the package root:
\begin{verbatim}
python -m pip install -r requirements.txt
python code/verify.py
python code/numerics.py
pdflatex -interaction=nonstopmode -halt-on-error article.tex
pdflatex -interaction=nonstopmode -halt-on-error article.tex
pdflatex -interaction=nonstopmode -halt-on-error article.tex
\end{verbatim}
The default exact run uses \pathtext{--rank-max 41 --even-max 12}.
Each coordinate in the JSON data is the tuple $(a,b,r,s)$ with the series
convention \eqref{eq:series}. The matrix keeps repeated nonzero rows. Kernel
columns are the $B_j$ family followed by the $C_j$ family. All matrix and
kernel entries are integers; scalar formulas use explicit rational
coefficients and named real atoms. The strings are data for the bundled
SymPy code, not assertions of numerical independence.

\pathtext{rank_certificates.json} records pivot-column information over the
fixed prime and the exact upper-bound checks.
\pathtext{weight5_certificate.json} contains the full central matrix and
witness. \pathtext{S_even_obstructions.json} records the uniform witness
through weight $41$. \pathtext{even_weight_identities.json} contains the
210 compiler outputs. \pathtext{leading_families.json} independently records
the 16 endpoint comparisons. \pathtext{verification_report.json} identifies
software versions and measured run time; that timing is environment-specific,
not an asymptotic performance claim. The numerical diagnostics have their
own file and expressly separate scope.

\begin{thebibliography}{9}
\raggedright
\bibitem{repo}
ProveIt Contributors,
\emph{Polylogarithms and their Arithmetic Bridges}, consolidated manuscript,
repository \texttt{VladimirReshetnikov/ProveIt}, commit
\pathtext{9bc738d3be22b8586a24693f19fb2e2a50ecd1bf},
particularly Chapter 5, the section containing
\pathtext{signed:conj:rank}, and Chapter 4's Gaussian product conventions.
\url{https://github.com/VladimirReshetnikov/ProveIt/tree/9bc738d3be22b8586a24693f19fb2e2a50ecd1bf/Analysis/Polylogarithms/docs/manuscript}.

\bibitem{zhao}
Jianqiang Zhao,
\emph{Standard relations of multiple polylogarithm values at roots of unity},
Documenta Mathematica \textbf{15} (2010), 1--34.
\url{https://doi.org/10.4171/DM/291}; arXiv:0707.1459.

\bibitem{panzer}
Erik Panzer,
\emph{The parity theorem for multiple polylogarithms},
Journal of Number Theory \textbf{172} (2017), 93--113.
Preprint arXiv:1512.04482v3 (2016).
\url{https://arxiv.org/abs/1512.04482}.

\bibitem{bachmann}
Henrik Bachmann and Khalef Yaddaden,
\emph{On a conjecture of Zhao related to standard relations among cyclotomic
multiple zeta values}, Journal of Algebra \textbf{688} (2026), 21--58.
Preprint arXiv:2411.18952.
\url{https://arxiv.org/abs/2411.18952}.

\bibitem{hirose}
Minoru Hirose,
\emph{Colored double zeta values and modular forms of general level},
arXiv:2205.08507.
\url{https://arxiv.org/abs/2205.08507}.
\end{thebibliography}
\end{document}
